\section{Integrated Worked Examples}
\label{sec:chapter1-integrated-examples}

The following examples combine several ideas from the chapter. Each solution
separates calculation from physical interpretation.

\begin{workedexample}
\textbf{Displacement from components.}

A particle moves $3\,\mathrm{m}$ east, $4\,\mathrm{m}$ north, and
$12\,\mathrm{m}$ upward. Find the displacement vector, its magnitude, and its
direction as a unit vector.

\begin{solution}
Using the Cartesian basis,
\[
\Delta\mathbf r=3\mathbf e_x+4\mathbf e_y+12\mathbf e_z\ \mathrm{m}.
\]
Its magnitude is
\[
\lVert\Delta\mathbf r\rVert
=\sqrt{3^2+4^2+12^2}\,\mathrm{m}=13\,\mathrm{m}.
\]
Therefore
\[
\widehat{\Delta\mathbf r}
=\frac{3}{13}\mathbf e_x+\frac{4}{13}\mathbf e_y
 +\frac{12}{13}\mathbf e_z.
\]
The unit vector contains only directional information; the factor
$13\,\mathrm{m}$ contains the size of the displacement.
\end{solution}
\end{workedexample}

\begin{workedexample}
\textbf{Work done by a constant force.}

A force $\mathbf F=(6,2,0)\,\mathrm{N}$ acts through the displacement
$\mathbf s=(3,-1,0)\,\mathrm{m}$. Calculate the work.

\begin{solution}
Work is the scalar projection of the force along the displacement multiplied by
the displacement magnitude, equivalently the dot product:
\[
W=\mathbf F\cdot\mathbf s
=(6)(3)+(2)(-1)=16\,\mathrm{J}.
\]
The positive result means that the force has a net component in the direction
of motion. A perpendicular component would contribute no work.
\end{solution}
\end{workedexample}

\begin{workedexample}
\textbf{Area and orientation from a cross product.}

Let $\mathbf a=(2,0,0)\,\mathrm{m}$ and
$\mathbf b=(1,3,0)\,\mathrm{m}$. Find the oriented area vector of the
parallelogram they span.

\begin{solution}
\[
\mathbf a\times\mathbf b
=\begin{vmatrix}
\mathbf e_x&\mathbf e_y&\mathbf e_z\\
2&0&0\\
1&3&0
\end{vmatrix}
=6\mathbf e_z\,\mathrm{m^2}.
\]
Hence the area is $6\,\mathrm{m^2}$ and the orientation is $+\mathbf e_z$ for
the ordering $\mathbf a\times\mathbf b$. Reversing the order reverses the
orientation but not the area.
\end{solution}
\end{workedexample}

\begin{workedexample}
\textbf{Decomposing a force parallel and perpendicular to a line.}

A force $\mathbf F=(5,4,0)\,\mathrm{N}$ acts near a rail directed along
$\mathbf d=(3,0,0)$. Find the components parallel and perpendicular to the rail.

\begin{solution}
The rail unit vector is $\widehat{\mathbf d}=\mathbf e_x$. Thus
\[
\mathbf F_{\parallel}
=(\mathbf F\cdot\widehat{\mathbf d})\widehat{\mathbf d}
=5\mathbf e_x\,\mathrm{N},
\]
and
\[
\mathbf F_{\perp}=\mathbf F-\mathbf F_{\parallel}
=4\mathbf e_y\,\mathrm{N}.
\]
The decomposition is verified by
$\mathbf F_{\parallel}\cdot\mathbf F_{\perp}=0$.
\end{solution}
\end{workedexample}

\begin{workedexample}
\textbf{Testing linear independence.}

Determine whether
$\mathbf v_1=(1,1,0)$, $\mathbf v_2=(0,1,1)$, and
$\mathbf v_3=(1,2,1)$ are linearly independent.

\begin{solution}
Direct inspection gives
\[
\mathbf v_3=\mathbf v_1+\mathbf v_2.
\]
Therefore
\[
\mathbf v_1+\mathbf v_2-\mathbf v_3=\mathbf 0
\]
is a nontrivial linear relation. The three vectors are linearly dependent and
span at most a two-dimensional subspace of $\mathbb R^3$.
\end{solution}
\end{workedexample}

\begin{workedexample}
\textbf{Angle between two vectors.}
For $\mathbf a=(1,1,0)$ and $\mathbf b=(1,0,1)$,
\[
\cos\theta=\frac{\mathbf a\cdot\mathbf b}{\lVert\mathbf a\rVert\lVert\mathbf b\rVert}
=\frac{1}{2},
\]
so $\theta=60^\circ$. The positive dot product signals an acute angle.
\end{workedexample}

\begin{workedexample}
\textbf{Parallel and perpendicular components.}
Resolve $\mathbf v=(4,2,0)$ relative to $\mathbf a=(1,1,0)$. Since
\[
\mathbf v_{\parallel}=\frac{\mathbf v\cdot\mathbf a}{\mathbf a\cdot\mathbf a}\mathbf a
=3(1,1,0)=(3,3,0),
\]
we obtain $\mathbf v_{\perp}=(1,-1,0)$. Directly,
$\mathbf v_{\perp}\cdot\mathbf a=0$.
\end{workedexample}

\begin{workedexample}
\textbf{Torque from a force.}
A force $\mathbf F=(0,12,0)\,\mathrm N$ acts at
$\mathbf r=(0.30,0,0)\,\mathrm m$. Then
\[
\boldsymbol\tau=\mathbf r\times\mathbf F=(0,0,3.6)\,\mathrm{N\,m}.
\]
The positive $z$ direction corresponds to counterclockwise rotational tendency in the $xy$ plane.
\end{workedexample}

\begin{workedexample}
\textbf{Area of a triangle.}
Let $P=(0,0,0)$, $Q=(3,0,0)$, and $R=(1,2,0)$. The area is
\[
\frac12\lVert(Q-P)\times(R-P)\rVert
=\frac12\lVert(3,0,0)\times(1,2,0)\rVert=3.
\]
The corresponding parallelogram has area $6$.
\end{workedexample}

\begin{workedexample}
\textbf{Testing linear independence.}
For columns $\mathbf v_1=(1,0,1)$, $\mathbf v_2=(0,1,1)$, and
$\mathbf v_3=(1,1,0)$,
\[
\det\begin{pmatrix}1&0&1\\0&1&1\\1&1&0\end{pmatrix}=-2\neq0.
\]
The vectors are therefore independent and form a basis of $\mathbb R^3$.
\end{workedexample}

\begin{workedexample}
\textbf{Coordinates in a rotated basis.}
Let $\mathbf e'_1=(\mathbf e_x+\mathbf e_y)/\sqrt2$ and
$\mathbf e'_2=(-\mathbf e_x+\mathbf e_y)/\sqrt2$. For
$\mathbf v=3\mathbf e_x+\mathbf e_y$,
\[
v'_1=\mathbf v\cdot\mathbf e'_1=2\sqrt2,
\qquad
v'_2=\mathbf v\cdot\mathbf e'_2=-\sqrt2.
\]
Both descriptions give $\lVert\mathbf v\rVert=\sqrt{10}$.
\end{workedexample}

\begin{workedexample}
\textbf{One Gram--Schmidt step.}
Take $\mathbf v_1=(1,1,0)$ and $\mathbf v_2=(1,0,1)$. Set
$\mathbf u_1=\mathbf v_1$ and compute
\[
\mathbf u_2=\mathbf v_2-
\frac{\mathbf v_2\cdot\mathbf u_1}{\mathbf u_1\cdot\mathbf u_1}\mathbf u_1
=\left(\frac12,-\frac12,1\right).
\]
After normalization,
$\widehat{\mathbf u}_1=(1,1,0)/\sqrt2$ and
$\widehat{\mathbf u}_2=(1,-1,2)/\sqrt6$.
\end{workedexample}

\begin{workedexample}
\textbf{Matrix acting on a vector.}
For
$T=\begin{pmatrix}2&1\\0&1\end{pmatrix}$ and $\mathbf v=(3,-1)$,
\[
T\mathbf v=(5,-1).
\]
The transformation stretches and shears the plane; it is not a rotation because it does not preserve all lengths and angles.
\end{workedexample}

\begin{workedexample}
\textbf{Eigenvectors of a diagonal matrix.}
For
$D=\operatorname{diag}(3,-2)$,
\[
D\mathbf e_x=3\mathbf e_x,
\qquad
D\mathbf e_y=-2\mathbf e_y.
\]
Thus the coordinate axes are eigenvector directions. The second direction is reversed because its eigenvalue is negative.
\end{workedexample}

\begin{workedexample}
\textbf{Functions as vectors.}
In the vector space of polynomials of degree at most two, write
$p(x)=2-3x+x^2$ in the basis $\{1,x,x^2\}$. Its coordinate vector is
$(2,-3,1)$. This example shows that a vector is an object admitting linear combination, not necessarily a geometric arrow.
\end{workedexample}
